% Author: Zhenyun Xie

\chapter*{Author's contributions} \label{sec:contributions}
\addcontentsline{toc}{chapter}{Author's contributions}
This is an industrial PhD thesis involving multiple disciplines, the research outcomes presented have been translated into two commercial products developed by iPronics: Smartlight and ONE.
The results presented in each chapter were achieved through extensive collaborative efforts spanning multiple disciplines, including photonic components design, PIC layout, fabrication, packaging, electronics, firmware, software development, etc.
The author's individual contributions to each chapter are summarized below:

\section*{Chapter 3}
For the applications results demonstrated on the Smartlight platform, I introduced the graph based methodology presented in Chapter 2 and implemented an automatic routing feature in the Smartlight SDK.
This feature was further applied to optimize the calibration process of the hexagonal PIC architecture.
At the application level, I developed software defined networking applications supporting arbitrary (N \times M) optical switching configurations.
In addition, I carried out the testing, validation and experimental characterization of the developed functionalities on the Smartlight product.

\section*{Chapter 4}
To address the scalability limitations of synthesizing optical circuit switching applications on Smartlight, I extended the graph based methodology to develop a software toolkit that supports multiple OCS architectures together with their corresponding routing algorithms.
I also co-designed and co-implemented the workflow for evaluating blocking performance.
Based on this development, I carried out blocking performance evaluations for different OCS architectures.
Combined with SNR simulation results generated using a complementary tool developed by the iPronics team, these studies contributed to the selection of the PIC architecture adopted for iPronics ONE.

\section*{Chapter 5}

Building upon the software development experience gained from Smartlight, I contributed to the development of the software stack for iPronics ONE.
My contributions include the implementation of the OCS simulation module, the definition of control interfaces, and the development of the software stack across multiple abstraction levels.
This work ranges from low-level control functions for individual photonic switching cells, to circuit-level control of the complete PIC, and finally to high-level connection management features.
On top of this framework, I co-developed the SDN control layer, enabling standardized control of iPronics ONE through the gNMI protocol.
